\section{Results}

We evaluated the proposed model family on a common anchored subset of the DTC test split and on external benchmarks designed to stress transfer beyond the source dataset. The overall pattern is consistent across settings: architectures that are strongest on same-dataset evaluation are not necessarily the ones that generalize best. Figure~\ref{fig:cross_dataset} summarizes this ranking shift across benchmarks, and Figure~\ref{fig:gen_gap} makes the same point more directly by plotting same-dataset performance against Davis transfer. The resulting generalization gap is the main empirical phenomenon of the paper.

\begin{figure}[t]
    \centering
    \includegraphics[width=\textwidth]{fig2_cross_dataset.png}
    \caption{Cross-dataset generalization across all benchmarks. All models are trained on DTC and evaluated on the same anchored interaction subsets. Anchor transfer variants (V2-650M, V2-35M) maintain strong performance under distribution shift, while pairwise baselines (DeepDTA, ESM-DTA) that excel on DTC collapse on external benchmarks.}
    \label{fig:cross_dataset}
\end{figure}

\begin{figure}[t]
    \centering
    \includegraphics[width=0.7\textwidth]{fig3_generalization_gap.png}
    \caption{Generalization gap: same-dataset AUROC (DTC test) versus cross-dataset AUROC (Davis). The dashed diagonal represents perfect generalization. Anchor transfer models (blue markers) cluster near the diagonal, while pairwise baselines (red squares) show large drops, indicating memorization of dataset-specific patterns rather than learning transferable binding features.}
    \label{fig:gen_gap}
\end{figure}

\subsection{Same-Dataset}

Table~\ref{tab:dtc_test} reports performance on the common DTC test subset ($n=27{,}536$). All sequence-based models are strong in this regime, with AUROC values between 0.870 and 0.910 (Figure~\ref{fig:same_dataset}). ESM-DTA is the best same-dataset model, reaching AUROC 0.910 together with the best CI and RMSE. Drug-Anchor and DeepDTA follow closely, while the two V2 variants lag slightly despite using the same frozen ESM-2 representation family.

\begin{figure}[t]
    \centering
    \includegraphics[width=0.85\textwidth]{fig5_same_dataset.png}
    \caption{Same-dataset AUROC on the DTC test set ($n=27{,}536$). All models perform well in-distribution, with ESM-DTA and Drug-Anchor leading. This ranking reverses under cross-dataset evaluation (cf.\ Figure~\ref{fig:cross_dataset}).}
    \label{fig:same_dataset}
\end{figure}

This same-dataset ranking is important because it establishes that stronger within-benchmark performance is not the same as stronger transfer. The best source-distribution model later becomes one of the weakest external models. Figure~\ref{fig:gen_gap} visualizes this reversal by contrasting DTC AUROC with Davis AUROC across methods.

\begin{table}[t]
\centering
\caption{Same-dataset performance on the common anchored DTC test set ($n=27{,}536$). Dashes indicate metrics unavailable for ConPlex*.}
\label{tab:dtc_test}
\small
\begin{tabular}{lccc}
\toprule
Model & AUROC & CI & RMSE \\
\midrule
ESM-DTA & \textbf{0.910} & \textbf{0.791} & \textbf{0.738} \\
Drug-Anchor & 0.906 & 0.786 & 0.783 \\
DeepDTA & 0.898 & 0.775 & 0.776 \\
V2-650M & 0.889 & 0.775 & 0.816 \\
V2-35M & 0.870 & 0.747 & 0.901 \\
ConPlex* & 0.768 & 0.680 & -- \\
\bottomrule
\end{tabular}
\end{table}

\subsection{Cross-Dataset Generalization}
\label{sec:cross_dataset}

The Davis benchmark produces the clearest cross-dataset ranking reversal. At test time, anchors are retrieved from the DTC training set by Tanimoto chemical similarity (chirality-aware Morgan fingerprints, radius 2, 2048 bits) after explicitly excluding all canonical chemical duplicates of Davis drugs from the retrieval pool. This ensures strict zero molecular overlap between training anchors and evaluation compounds. We verified overlap at six levels (Table~\ref{tab:overlap_audit}): while raw SMILES string overlap is 0\%, canonical SMILES overlap is 83.8\% and InChIKey overlap is 89.7\%, confirming that most Davis kinase inhibitors are present in DTC under different string representations. All such duplicates are removed before anchor retrieval.

\begin{table}[t]
\centering
\caption{Drug overlap audit between Davis and the DTC training set at six levels of chemical identity. Raw SMILES string comparison substantially underestimates true overlap.}
\label{tab:overlap_audit}
\small
\begin{tabular}{lcc}
\toprule
Identity level & Overlap & \% \\
\midrule
Raw SMILES string & 0/68 & 0.0 \\
Canonical SMILES & 57/68 & 83.8 \\
Full InChIKey & 61/68 & 89.7 \\
First-block InChIKey (no stereo) & 63/68 & 92.6 \\
Morgan FP (no chirality) & 59/68 & 86.8 \\
Morgan FP (with chirality) & 57/68 & 83.8 \\
\bottomrule
\end{tabular}
\end{table}

After duplicate exclusion, the retrieval achieves 99.7\% coverage ($n=29{,}975$). Table~\ref{tab:davis_cross_dataset} reports results with AUPRC to account for the severe class imbalance (8.6\% binders). V2-650M leads with CI 0.642, AUROC 0.669, and AUPRC 0.209. The absolute drops from Table~\ref{tab:dtc_test} are substantial for pairwise models: DeepDTA and ESM-DTA collapse to near-random. Figure~\ref{fig:gen_gap} highlights this as a generalization-gap effect.

Table~\ref{tab:tanimoto_strat} stratifies V2-650M performance by Tanimoto similarity between the query drug and the retrieved anchor's drug, directly testing how the model performs when the anchor's chemistry differs from the query.

\begin{table}[t]
\centering
\caption{Cross-dataset generalization from DTC to Davis ($n=29{,}975$) after excluding canonical chemical duplicates. Anchors retrieved by Tanimoto similarity from the DTC training set. Per-protein metrics.}
\label{tab:davis_cross_dataset}
\small
\begin{tabular}{lcccc}
\toprule
Model & CI $\uparrow$ & AUROC $\uparrow$ & AUPRC $\uparrow$ & RMSE $\downarrow$ \\
\midrule
V2-650M & \textbf{0.642} {\scriptsize [.634--.651]} & \textbf{0.669} {\scriptsize [.653--.686]} & \textbf{0.209} {\scriptsize [.192--.227]} & \textbf{0.923} \\
V2-35M  & 0.591 {\scriptsize [.583--.599]} & 0.618 {\scriptsize [.603--.634]} & 0.171 {\scriptsize [.157--.186]} & 1.069 \\
DeepDTA & 0.521 {\scriptsize [.513--.528]} & 0.504 {\scriptsize [.487--.521]} & 0.155 {\scriptsize [.141--.170]} & 1.157 \\
ESM-DTA & 0.501 {\scriptsize [.495--.507]} & 0.503 {\scriptsize [.488--.519]} & 0.140 {\scriptsize [.129--.152]} & 1.299 \\
\bottomrule
\end{tabular}
\end{table}

\begin{table}[t]
\centering
\caption{V2-650M performance stratified by Tanimoto similarity between the query drug and the retrieved anchor's drug. After excluding canonical duplicates, no exact matches remain. The model generalizes even to chemically dissimilar compounds (Tanimoto~$< 0.6$), with the best performance at moderate similarity (0.6--0.8).}
\label{tab:tanimoto_strat}
\small
\begin{tabular}{lccccc}
\toprule
Tanimoto bin & $n$ & Drugs & CI $\uparrow$ & AUROC $\uparrow$ & AUPRC $\uparrow$ \\
\midrule
{[0--0.6)}   & 14,992 & 34 & 0.656 & 0.651 & 0.224 \\
{[0.6--0.8)} & 7,938  & 18 & \textbf{0.708} & \textbf{0.778} & \textbf{0.627} \\
{[0.8--0.95)} & 7,045 & 16 & 0.539 & 0.539 & 0.309 \\
\midrule
Overall      & 29,975 & 68 & 0.642 & 0.669 & 0.209 \\
\bottomrule
\end{tabular}
\end{table}

The [0.6--0.8) bin achieves the best performance (CI 0.708, AUROC 0.778), while the [0.8--0.95) bin drops to CI 0.539. This non-monotonicity likely reflects the small sample size (16 drugs in the high-similarity bin) rather than a systematic failure mode; the confidence intervals across bins overlap substantially.

\subsection{Benchmark-Wise Anchor Quartile and Chemical Novelty Analysis}
\label{sec:benchmark_heatmaps}

The Davis and Metz benchmarks answer noticeably different questions before and after overlap control, so Figures~\ref{fig:benchmark_heatmaps_main} and~\ref{fig:benchmark_quartile_dists_main} report both the unfiltered and filtered protocols. The unfiltered view shows nominal benchmark difficulty after only model-coverage checks, whereas the filtered view removes proteins and drugs overlapping the DTC training reference before benchmark-internal anchors are defined. We compare V2 in oracle, weakest-anchor, and random-anchor modes against the paper's sequence/SMILES baselines (DeepDTA, ConPlex*, and ESM-DTA). DrugBAN is omitted here because its graph-drug encoder changes the molecular input modality and is therefore not a like-for-like comparison to the main baseline family.

The paired views are informative. On unfiltered Davis ($n=30{,}056$, 442 proteins, 68 drugs), V2-oracle achieves mean per-protein CI 0.605, ahead of DeepDTA (0.521) and ESM-DTA (0.508), while ConPlex* remains competitive at 0.573. After overlap filtering, Davis contracts sharply to $n=854$ interactions over 122 proteins and only 7 drugs; in this stricter regime V2-oracle rises to mean per-protein CI 0.631, compared with 0.599 for DeepDTA and 0.476 for ESM-DTA, while ConPlex* reaches 0.637. Metz shows the opposite pattern: in the unfiltered benchmark ($n=35{,}259$, 170 proteins, 1,423 drugs), all methods cluster tightly, with V2-oracle at 0.614 and DeepDTA at 0.609, indicating that nominal Metz is heavily overlap-dominated. Once Metz is filtered to $n=374$ interactions over 9 proteins and 115 drugs, V2-oracle becomes the clearest winner at mean per-protein CI 0.644, above DeepDTA (0.583), ConPlex* (0.586), and ESM-DTA (0.490).

\begin{figure}[t]
    \centering
    \begin{subfigure}[t]{0.49\textwidth}
        \centering
        \includegraphics[width=\textwidth]{fig_davis_unfiltered_heatmap_ci.png}
        \caption{Davis, unfiltered}
    \end{subfigure}\hfill
    \begin{subfigure}[t]{0.49\textwidth}
        \centering
        \includegraphics[width=\textwidth]{fig_davis_filtered_heatmap_ci.png}
        \caption{Davis, filtered}
    \end{subfigure}

    \vspace{0.8em}

    \begin{subfigure}[t]{0.49\textwidth}
        \centering
        \includegraphics[width=\textwidth]{fig_metz_unfiltered_heatmap_ci.png}
        \caption{Metz, unfiltered}
    \end{subfigure}\hfill
    \begin{subfigure}[t]{0.49\textwidth}
        \centering
        \includegraphics[width=\textwidth]{fig_metz_filtered_heatmap_ci.png}
        \caption{Metz, filtered}
    \end{subfigure}
    \caption{Per-protein mean CI heatmaps stratified jointly by anchor pKi quartile and maximum Tanimoto similarity to the DTC training drug library. The top row shows Davis before and after overlap filtering; the bottom row shows the same comparison for Metz. Each panel compares V2 oracle/weakest/random controls against DeepDTA, ConPlex*, and ESM-DTA. Numbers in cells indicate mean per-protein CI with the number of interactions in parentheses. The unfiltered panels reflect nominal benchmark difficulty, whereas the filtered panels expose how strongly the conclusions depend on overlap control.}
    \label{fig:benchmark_heatmaps_main}
\end{figure}

\begin{figure}[t]
    \centering
    \begin{subfigure}[t]{0.49\textwidth}
        \centering
        \includegraphics[width=\textwidth]{fig_davis_unfiltered_quartile_ci_distribution.png}
        \caption{Davis, unfiltered}
    \end{subfigure}\hfill
    \begin{subfigure}[t]{0.49\textwidth}
        \centering
        \includegraphics[width=\textwidth]{fig_davis_filtered_quartile_ci_distribution.png}
        \caption{Davis, filtered}
    \end{subfigure}

    \vspace{0.8em}

    \begin{subfigure}[t]{0.49\textwidth}
        \centering
        \includegraphics[width=\textwidth]{fig_metz_unfiltered_quartile_ci_distribution.png}
        \caption{Metz, unfiltered}
    \end{subfigure}\hfill
    \begin{subfigure}[t]{0.49\textwidth}
        \centering
        \includegraphics[width=\textwidth]{fig_metz_filtered_quartile_ci_distribution.png}
        \caption{Metz, filtered}
    \end{subfigure}
    \caption{Quartile-wise per-protein CI distributions under the same paired protocols. The unfiltered panels show the benchmark as typically reported, while the filtered panels remove protein and chemical overlap with the DTC training reference. This makes the regime change visible directly in the distributions: unfiltered Metz compresses the methods into similar CI ranges, whereas filtered Metz widens the gap in favor of V2-oracle; Davis remains more separated across anchor quartiles in both protocols, but the filtered view is much smaller and chemically narrower.}
    \label{fig:benchmark_quartile_dists_main}
\end{figure}

\subsection{Ablation: ESM-2 Scaling and Baseline Controls}

The ablations reveal that the transfer mechanism is sensitive to protein encoder scale (Figure~\ref{fig:esm_scaling}). Scaling from ESM-2 35M to 650M improves Davis performance (AUROC 0.669 versus 0.644, CI 0.642 versus 0.605). The larger model captures richer protein representations that improve the anchor--query comparison.

\begin{figure}[t]
    \centering
    \includegraphics[width=0.75\textwidth]{fig7_esm_scaling.png}
    \caption{Effect of ESM-2 scaling (35M vs 650M parameters) on Davis cross-dataset transfer. The larger model improves performance across all anchor quartiles, with the largest gains on weak anchors (Q1--Q2).}
    \label{fig:esm_scaling}
\end{figure}

ESM-DTA serves as a negative control for the hypothesis that stronger protein embeddings alone are enough. It achieves the best same-dataset performance on DTC (AUROC 0.910, CI 0.707) but collapses on Davis (per-protein AUROC 0.503, CI 0.501---indistinguishable from random). The contrast between Table~\ref{tab:dtc_test} and Table~\ref{tab:davis_cross_dataset} indicates that the improvement does not come from ESM-2 alone. It comes from changing the prediction problem from isolated pair scoring to transfer from a known anchor.

\subsection{Secondary Controls: Drug-Side Anchors and Conformation Weighting}

Two secondary analyses help delimit what the anchor mechanism is, and is not, using as signal. First, the Drug-Anchor ablation reverses the modality of the support example by conditioning on a strong \emph{drug} for the same protein rather than a strong protein for the same drug. This variant remains competitive in-distribution (DTC test AUROC 0.906) and stays above pairwise baselines on several external benchmarks, but it does not match the main protein-anchor formulation on Davis (AUROC 0.559 in Figure~\ref{fig:cross_dataset}, versus 0.669 for V2-650M). This makes the ablation useful as a control rather than a replacement: the transfer idea is not protein-exclusive, but protein-side anchoring is the stronger direction for the kinase-focused cross-dataset setting studied here.

Second, a conformation-weighted negative control on the IDP benchmark asks whether structural proximity alone can explain the gain. Replacing the query ESM-2 embedding with a TM-score-weighted average of embeddings from structurally matched ordered proteins reduces AUROC from 0.584 to 0.575 on the common IDP subset ($n=2{,}442$). The result argues against a simple ``borrow a similar structure'' explanation. What transfers is not generic structural resemblance, but experimentally grounded binding context supplied by the anchor. Full architectural and protocol details for both controls are provided in the Supplementary Information.
